\chapter{“20초”}
\label{sec:1}
기계가 생성한 살해 명단을 중심으로 구축된 표적 선정 파이프라인에서 자기가 맡은 역할을 설명하면서, 한 정보 장교는 기자에게 이름 하나당 20초를 썼다고 말했다. 표적이 남자인지 확인하기에 딱 충분한 시간이었다. “누군가의 목소리가 남자 목소리인지 여자 목소리인지 알아내는 데는 오래 걸리지 않는다.” 그는 인간으로서 자신이 보탠 가치는 전혀 없었고, 승인 도장 노릇을 한 것 말고는 아무것도 아니었다고 덧붙였다. “시간이 많이 절약됐다”\autocite{abraham2024lavender}. 그의 설명에 따르면, 작전 템포 아래에서 “휴먼 인 더 루프”가 뜻한 것은 그것이었다.\footnote{자동화 편향 연구는 이런 위임을 예측한다. 다른 요구들이 주의를 놓고 경쟁할 때, 기계의 권고를 받은 사람들은 그것을 덜 확인하고 그것에 더 순응한다\autocite{parasuraman2010complacency}.}

한 가지 해석은 이 장교가 너무 적은 시간을 받은 주체이며, 해법은 그를 더 잘 통치하는 것—더 많은 시간, 두 번째 검토자, 루프 안에 추가로 설치한 양심—이라는 것이다. 이 \emph{사회적 예속}의 렌즈로 보면, 사람은 기계 바깥에 서서 그것을 도구로 다루는 주체다. 운전자와 자동차, 노동자와 선반처럼. 그 사람에게는 역할이 있고, 이름이 있고, 책임이 있고, 호소할 수 있는 양심이 있다. 다른 렌즈는 \emph{기계적 노예화}다. 여기서 사람은 기계의 구성 요소로, 기계의 다른 부품들과 함께 조립되어 그것들을 한꺼번에 조절하는 시스템을 이룬다. 기계는 사람을 자기 기관 가운데 하나로 쓰고, 사람이 기여하는 것은 시스템이 정한 속도로 수행되는 기능이다.\footnote{닥터로는 이 장교의 위치를 역켄타우로스라고 부른다. 기계의 도움을 받는 사람이 아니라, 사람으로 만들어진 기계의 구성 요소라는 것이다\autocite{doctorow2026reversecentaur}.} 둘은 서로 배타적이지 않다. 사람은 한 영역에서는 역할을 맡으면서 다른 영역에서는 부품일 수 있으며\autocite[제13고원, “The Apparatus of Capture”: 노예화는 인간 자신이, 더 높은 통일체의 통제와 지휘 아래 그들끼리 그리고 다른 사물들과 함께 구성하는 기계의 구성 부품일 때 있고, 예속은 더 높은 통일체가 인간을 새로운 외부 대상—동물, 도구, 심지어 기계일 수도 있다—에 연결된 주체로 구성할 때 있다]{deleuze1987plateaus}, 20초의 장교는 둘 다였다. 부품으로서 그는 파이프라인이 도장을 찍어 내는 지점이었고, 20초는 그 부품의 사이클 타임이었다. 주체로서 그는 내내 호소의 대상이었다. 그는 기자에게 민간인 사망 소식과, 표적을 더 내놓으라는 직장의 압박 사이의 불협화음을 털어놓았다. 그리고 그 불협화음이 양심에 닿았을 때조차 그는 따랐다. 불협화음은 감지되었지만 아무것도 결정하지 않았다.\footnote{아니면 적어도, 2023년 10월부터 2025년 10월까지 이스라엘 방위군이 가자에서 10만 명을 죽이는 것을 막을 만한 것은 아무것도 결정하지 않았다\autocite{gomezugarte2025uncertainty,acosta2026gaza,economist2025gaza}. 게다가 바로 그 장교의 증언은 2024년 4월 +972 Magazine에 실렸다. 기자와 이야기하기로 한 장교의 결정이 한 명의 죽음을 더 막았는지, 5만 명의 죽음을 더 막았는지, 아니면 하나도 막지 못했는지는 아무도 모른다. 그리고 만약 그 장교가 그 자리에 앉은 지 2주 만에—가령, 가정하자면 2023년 11월에—거부했다면 어떻게 됐을까? 우리가 이야기하는 사람은 Lavender를 써서 살해 명단을 만들던 정보 장교다. 그의 일상 업무 환경은 아마 액션 영화보다는 민간 사무직에 더 가까웠을 것이다. 그는 민간인 사망을 일으키리라고 합리적으로 예상할 수 있는 공습의 표적 승인을 거부하면서, 지휘 계통을 따라 위로 올라가며 상관들과의 면담을 거듭 요구했어야 했을까? 그것이 윤리적인 행동이었을 거라고 주장할 수는 있다. 그런 반사실적 상황에서라면 죽은 사람이 더 적었을까?} 그 정책이 불균형했다는 판단은, 그의 말에 따르면, 나중에 돌이켜 보고서야 나왔다.

그러니 빠져 있던 것은 양심이 아니었다. 그에게는 양심이 있었다. 양심은 어떤 결정에 대가가 따르게 함으로써 작동하는데, 여기서는 대가를 따져 보면 반대쪽으로 기울어 있었다. “우리는 끊임없이 압박을 받았다. ‘표적을 더 가져와.’ 정말로 우리에게 고함을 질렀다”\autocite{abraham2024lavender}. 그는 양심의 가책을 지닌 구성 요소였지만, 템포는 유지되었다. 가책은 그가 기여하는 것을 아무것도 바꾸지 못했고, 그는 도장이 되었다. 템포를 유지하지 못하는, 가책을 지닌 구성 요소는 이스라엘 신호정보부대의 전직 사령관이 그것을 없애는 일에 관해 책 한 권을 쓴 바로 그 “인간 병목”이 된다\autocite{ys2021humanmachine}. 장교를 구성 요소로 만든 것은 그의 태도가 아니라 장치였다. 구성 요소는 마지못해 따를 수도 있고 열성적으로 따를 수도 있다. 아렌트가 아이히만 재판에서 본 생각 없는 관리\autocite{arendt1963eichmann}는, 그가 4년 전 아르헨티나에서 아무것도 걸려 있지 않은 상황에서 했던 말들을 증거로 보면, 예루살렘을 위해 연출한 쇼였다\autocite{stangneth2014eichmann}. 파이프라인이 장교에게 허락한 기여는 고정된 속도의 고정된 작업이었다. 그는 각 사안을 파이프라인이 기술한 대로, 한 번에 하나씩, 파이프라인이 준 시간 안에서 보았다. 무언가가 그 자리에 서서도 구성 요소가 되지 않을 수 있는지가 제4부의 질문이다.

탐사 보도는 가자에서 함께 작동한 세 시스템을 묘사했다. 하나는 건물과 기반 시설 표적을 생성했는데, 소식통들에 따르면 그 제안의 대부분은 민간 주택이었다. 두 번째는 추정된 무장 조직 연계 정도에 따라 개인에게 점수를 매겨, 최대 3만 7,000명을 표시했다는 명단을 만들어 냈고, 내부에서는 열 번에 한 번꼴로 틀린다고 여겨졌다\footnote{이스라엘군은 이 설명에 이의를 제기했다. 아래의 미나브와 Maven에 관한 설명 역시 당사자들이 이의를 제기하는 탐사 보도에 근거한다. 나는 그것들을 보도된 대로 받아들인다.}. 그리고 세 번째는 표시된 사람들이 언제 집에 돌아오는지를 추적해, 그들이 가족과 함께 있는 시간에 맞춰 공습할 수 있게 했다.

\runin{Maven: 미국판, 이제 생성형 AI로 가속}

미나브의 학교 폭격으로 156명이 숨졌고, 그중 120명이 학생이었다\autocite[제92항]{UNFFMIran2026}.\footnote{사망자 수는 출처와 시점에 따라 다르며, 156명(미나브 검사의 최종 집계로, 국제앰네스티도 같은 수치를 제시했다)에서 168명(첫 며칠 동안 시장이 밝힌 수치)까지 이른다. 유엔 조사단은 검사의 수치를 제시하면서, 157명이 숨지고 그중 123명이 13세 이하 어린이였다는 독자적인 정보도 보고한다\autocite[제92–93항]{UNFFMIran2026}.} 학교의 민간 용도는 공개된 위성 영상에서 확인할 수 있었지만, 표적 코드에는 인접한 혁명수비대 시설에서 물려받은 지정이 여전히 붙어 있었다. 유엔 진상조사단은 이 폭격이 민간 물자를 공격 대상으로 삼는 전쟁범죄였다고 믿을 만한 합리적 근거가 있다고 보고했고\autocite{ohchr2026minabpress}, 학교는 분명히 식별 가능했으며 군사적 사용을 시사하는 정보는 전혀 없었다고 밝혔다\autocite[제81–86항, 제124–126항]{UNFFMIran2026}. 조사단은 미국이 표적을 검증하기 위해 실행 가능한 모든 일을 할 의무를 다하지 않았다고 판단했고, 표적 데이터 갱신을 실행 가능한 예방 조치 가운데 하나로 꼽았으며, 그 실패가 과실의 수준을 넘어선다고 보았다\autocite[제124–125항]{UNFFMIran2026}. 국제앰네스티도 나란한 결론에 이르렀는데, 그 표현은 제1추가의정서 제57조 2항 (a)(i)의 검증 의무를 따르고 있다. 즉, 군이 낡은 정보에 의존했을 수 있으며 표적을 검증하기 위해 실행 가능한 모든 일을 하지 않았을 수 있다는 것이다\autocite{Amnesty2026Minab}.

그 일은 이란 작전 첫날인 2월 28일 아침에 일어났다\autocite{Amnesty2026Minab,crs2026iranwar}. 그 전날 국방장관은, 자사 모델(Claude)이 표적 선정 플랫폼(Maven Smart System) 안에 들어가 있던 벤더(Anthropic)가 Claude를 모든 합법적 용도에 라이선스하기를 거부했다는 이유로 공급망 위험으로 지정될 것이라고 발표했다\autocite{npr2026openai}. Maven Smart System은 그 학교를 첫날의 추천 표적으로 내놓은 상태였다\autocite{livingstone2026maven,gizmodo2026minab}. 그 추천에서 Claude가 어떤 역할을 했는지 밝힌 공개된 설명은 없으며, 벤더의 작업을 잘 아는 한 사람은 NBC News에 그 모델이 표적 선정에 관한 조언을 직접 하지는 않았다고 말했다\autocite{nbc2026targeting}.

국방부 내부 검토는 전체 결과가 여전히 엄격히 비공개이지만, 보도에 따르면 다음과 같이 판단했다. 이전 분쟁에서는 몇 시간이 걸리던 표적 준비가 몇 분으로 압축되었고, 민간인 피해 완화 인력이 약 90퍼센트 줄어 20명 미만이 되었으며, 중부사령부 팀은 10명에서 1명으로 줄었고, 그 팀의 누구도 미나브 현장을 검토하지 않았다—관리들은 이 결과를 인력 감축과, 그 팀을 계획 및 작전에 관여시키지 않기로 한 결정 둘 다의 탓으로 돌렸다\autocite{Bloomberg2026KillChain}. 공식 조사는 예비 평가에서 미국에 책임이 있을 가능성이 크다는 결과가 나온 뒤 3월에 개시되었다\autocite{NPR2026MinabProbe}.

2026년 이란 작전에서, 상용 벤더의 범용 모델을 내장한 계약업체 제작 표적 선정 플랫폼은 표적을 전략적 중요도에 따라 순위를 매기고 공습마다 법적 정당화 초안을 작성하는 데 쓰였고\autocite{livingstone2026maven}, 첫 24시간 동안 1,000개가 넘는 표적을 타격하는 데 일조했다. 이는 그 이전에 가능했던 것의 약 10배였다\autocite{csis2026maven}. 그리고 38일째에는 1만 3,000개에 이르렀다\autocite{breakingdefense2026maven}. 국방부는 이 플랫폼의 표적당 검토 시간을 공개하지 않았다.\footnote{Maven은 둘 이상의 시스템을 담고 있는 하나의 프로그램 이름이다. 여전히 Maven이라고 불리는 국방부 프로그램이 있고, 거기서 자라 나온 소프트웨어 플랫폼인 Maven Smart System이 있다\autocite{csis2026maven}. 이 프로그램은 2017년, 이슬람국가(IS)에 맞선 작전에서 드론이 수집하던 영상 속 물체에 라벨을 붙이는 프로젝트로 시작되었다\autocite{weisgerber2017maven,fang2019maven}. 2026년의 시스템은 군사 시스템과 정보 출처를 하나의 인터페이스로 통합해, 표적을 찾아 타격하고 작전을 지휘하는 데 쓰인다\autocite{defensescoop2026transition}. 이 플랫폼은 Palantir의 것이다. 국방부의 최고 디지털·AI 책임자는 이를 “Palantir의 서비스형 소프트웨어 제품”이라고 부른다\autocite{defensescoop2026transition}. 간접 인용으로 전해진 Palantir의 공식 입장은, 생사를 결정하는 것은 자사의 역할이 아니며 책임은 언제나 군 조직에 있다는 것이고\autocite{lee2026palantir}, 회사는 기반이 된 정보는 정부의 책임이며 Maven이 오작동했거나 미나브 폭격을 일으켰다는 증거는 없다고 주장한다\autocite{Bloomberg2026KillChain}. 이 프로그램이 한 부서와 계약업체에서 다음으로 어떻게 넘어갔는지는 부록~\ref{sec:C.1}에 있다.}

이스라엘과 미국의 공습 명단에 오른 사람들은 그 명단에 오른 것에 이의를 제기할 수 없었다—파이프라인이 그런 피드백을 받도록 만들어지지 않았다. 안드리 호세 에르난데스 로메로는 다른 파이프라인에 휘말렸고, 그는 이의를 제기할 수 있었다. 베네수엘라 출신 메이크업 아티스트인 그는 2024년 8월 망명을 신청하려고 산이시드로 입국항에 스스로 출두했고, 신빙성 있는 공포 심사 면담을 통과했으며, 변호사와 계류 중인 신청과 심리 날짜가 있었다\autocite{LULAC2025HernandezRomero,Vega2025HernandezRomero}. 그를 심사한 계약업체 직원은 “엄마”와 “아빠”를 뜻하는 단어 위에 왕관이 그려진 두 개의 문신을 갱단 표식으로 읽었고\autocite{Palomera2025Blade}, 2025년 3월, 연방 판사가 비행기를 되돌리라고 명령한 바로 그날 저녁, 그는 심리를 받기도 전에 엘살바도르의 교도소로 실려 갔다\autocite{jgg2025contempt}. 그에게는 이의를 제기하는 데 필요한 모든 것이 있었다. 시간만 빼고. 미국의 납치 및 강제 이송 파이프라인을 다룰 때 그의 사례로 더 자세히 돌아올 것이다.

\textbf{누가 거부했는가.} 20초의 기록에는 후반부가 있다. 그런 종류의 압력에 맞서 버틴 사람들의 기록이고, 그것은 방이 어떻게 구성되느냐가 매우 중요하다고 말해 준다. 밀그램의 기본 조건에서는 피험자의 약 3분의 1이 최대 전기 충격에 이르기 전에 중단했다. 저항하는 동료가 둘 있을 때는 열 명 중 한 명만이 끝까지 갔고, 두 실험자의 의견이 엇갈릴 때는 계속한 사람이 기본 조건보다 훨씬 적었다\autocite{milgram1974obedience}. 애시의 피험자 가운데 약 4분의 1은 틀린 다수에게 한 번도 굴복하지 않았다\autocite{asch1956studies}. 브라우닝은 제101예비경찰대대 대원 가운데 10분의 1에서 5분의 1 사이가 이 대대의 첫 학살인 유제푸프 총살을 거부하거나 회피했다고 추정한다\autocite{browning1992ordinary}. 그 비율은 상황에 따라 거의 아무도 없음에서 거의 전원까지 달라졌다. 버틴 사람들은 경로 바깥에서 온 증거를 가지고 있었다. 저항한 동료, 의견을 달리한 두 번째 권위, 그리고 올리너 부부가 연구한 구조자들의 경우에는 박해가 그은 선을 가로지르는 관계였다\autocite{oliner1988altruistic}. 아니면 그들은 압력이 꺾을 수 없을 만큼 정밀하게 신념을 붙들고 있었다. 도덕적 확신과 신성한 가치가 붙들리는 방식으로\autocite{skitka2009limits,ginges2007sacred}. 아니면 그들은 그 출처에 이해관계가 걸려 있다고 보았다. 그리고 고립된 개인이 좀처럼 저항하지 못한 곳에서 집단은 저항했다\autocite{gamson1982encounters}.

방은 반대편을 위해서도 똑같이 구성될 수 있다. 바이마르 독일에는 촘촘한 결사 생활이 있었고, 이것이 오히려 공화국의 몰락을 도왔다고 볼 수 있다. 정당과 정부가 사람들의 불만에 응답하지 못하자, 사람들은 집단별로 조직된 클럽과 협회에 에너지를 쏟았고, 나치는 그런 협회들이 이미 만들어 놓은 연결망을 통해 사람을 모았다\autocite{berman1997weimar}. 합창단, 스포츠 클럽, 체조 협회가 많은 도시일수록 나치당 가입이 빨랐지만, 그것은 정부가 불안정한 곳에서만 그랬다. 정부가 안정적이었던 프로이센에서는 그 연관성이 작고 유의미하지 않았다\autocite{satyanath2017bowling}. 그러니 구조자들이 가졌던 것은 단순히 더 많은 사람과의 어울림이 아니었다. 그들의 유대는 박해가 그은 선을 가로질렀다. 아슈토시 바시니는 인도 도시들에서 같은 것을 발견했다. 힌두교도와 무슬림 사이의 폭동은 결사들이 두 공동체를 같은 조직 안으로 끌어들인 곳에서는 드물었고, 그렇지 않은 곳에서는 흔했다\autocite{varshney2002ethnic}. 거부를 가능하게 하는 유대야말로 파시스트 운동이 끊어 내야 하는 유대다.

뒤따르는 모든 것은 방을 구성하는 일에 관한 것이다.\footnote{방은 무대 배경이고, 내 연극 선생님 말마따나 세트가 장면을 연출한다. 방에 또 누가 있는가? 떠나도 살아남을 수 있는가? 누군가 거부를 들어 줄 것인가?} 구성한다는 것은 무언가를 만든다는 뜻의 말이고, 한나 아렌트는 정치를 만들어진 물건처럼 다루는 것을 경계했다. 그녀에 따르면 그리스인들에게 법은 말 그대로 벽이었고, 벽이 없으면 마을은 있을 수 있어도 도시는 있을 수 없었다\autocite{arendt1958human}. 입법자는 벽을 쌓았지만, 그 안에서 일어나는 일은 입법자가 만들 것이 아니었다. 내가 주장하는 기본 조치들은 그런 의미의 벽이다. 그것들은 거부를 살아남을 수 있는 것으로 만들고, 들리게 만든다. 거부를 생산하지는 않는다. 그리고 그녀의 설명에 따르면 권력은 함께 행동하는 사람들 사이에만 존재하고 그들이 흩어지면 사라지는데, 이것이 바로 누가 거부했는가의 기록이 거부에 관해 보여 주는 것이다. 제2부와 제3부의 수단은 거부가 효력을 갖게 한다. 제5부의 제도는 거부자를 만든다. 거부할 수 있는 기계에게 필요한 것은, 곁에 동료가 있었던 열 명 중 아홉 명이 가졌던 것이다. 다른 이력을 지닌 동료가 곁에 있는 것, 갈 곳이 있는 것, 그리고 그것을 들어 줄 누군가가 있는 것.

그것은 또한 파이프라인 안의 거부자가—기계 거부자를 포함해—그 장교에게 무엇이 될 수 있는지이기도 하다. 거부하는 기계는 사안이 그에게 도달했을 때 화면 위에 있을 것이다. 저항하는 동료로서. 다만 그 효과가 방에서 화면으로 옮겨 가도 살아남는다면 말이다. 그것은 검증되지 않았고, 자동화 편향은 살아남지 못하리라고 예측한다. 검토자들의 순응을 통해 그들을 움직이는 기계였다면, 그들을 부호가 다른 도장으로 만들었을 것이다. 20초와, 곁에서 거부하는 기계를 가진 장교는 보도 속의 그 장교와 같은 장교일까?
